\documentclass[../DoAn.tex]{subfiles}
\begin{document}

\section{Đặt vấn đề}
\label{section:1.1}

Giao thông đô thị Hà Nội là dòng xe hỗn hợp, trong đó xe máy chiếm đa số. Xe máy chiếm trên 90\% số phương tiện đăng ký tại Hà Nội~\cite{nguyen2022motorcycles}. Tại nút Láng Hạ -- Láng, một nút có đèn trên cùng trục Láng Hạ -- Lê Văn Lương, xe máy chiếm khoảng 86,5\% số phương tiện đếm từ video trong giờ cao điểm~\cite{mai2019signal}. Thành phố Hồ Chí Minh có tỷ lệ tương tự, 85--90\%~\cite{nguyen2026surtrac}. Ô tô và xe máy dùng chung lòng đường; xe máy không đi theo làn cố định mà len vào các khoảng trống ngang giữa các xe khác.

Dòng xe tại nút có đèn còn có nhiều hành vi trái luật. Video tại hai nút ở Hà Nội cho thấy 9,64--11,06\% xe máy vượt đèn đỏ và 13,55--15,82\% xe máy xuất phát sớm, tức bắt đầu đi từ khoảng 5~s trở lên trước khi đèn xanh. Với ô tô, các tỷ lệ này chỉ từ 0,50\% đến 1,61\%~\cite{mai2019signal}. Cùng nghiên cứu ghi nhận một lượng lớn xe máy của luồng phụ đi vào giữa nút và dừng chờ khi pha đèn đổi~\cite{mai2019signal}. Luận văn xem xét 11 hành vi như vậy; danh sách đầy đủ ở Bảng~\ref{table:1.1}.

SUMO~\cite{lopez2018sumo} là trình mô phỏng vi mô phổ biến để đánh giá phương án tổ chức giao thông trước khi triển khai. Mô hình mặc định của SUMO giả định xe bám theo xe phía trước trong cùng làn, dừng trước vạch khi đèn đỏ và chỉ vào nút khi được phép. Mô hình làn con (sublane) cho phép nhiều xe hẹp chiếm chung một làn và di chuyển theo phương ngang. Dù vậy, với tham số mặc định, SUMO không tạo ra vượt đèn đỏ có xác suất, xuất phát sớm, vượt đèn theo đoàn hay dừng chờ có chủ đích giữa nút. Một mô hình SUMO mặc định cho nút giao ở Hà Nội vì thế có thể cho hàng chờ, thời gian thông qua và số xung đột khác xa thực tế.

Để tái hiện các hành vi trên, người lập mô hình phải chỉnh cùng lúc ba nhóm tham số và một nhóm luật:
\begin{itemize}
    \item tham số bám xe: thời gian phản ứng, khoảng cách tối thiểu, gia tốc, giảm tốc;
    \item tham số chuyển làn và di chuyển ngang của mô hình làn con;
    \item tham số ứng xử tại nút: thời gian chấp nhận đi tiếp khi đèn đỏ, khoảng trống chấp nhận;
    \item các luật có điều kiện mà tham số tĩnh không biểu diễn được.
\end{itemize}
Không gian tham số này lớn, các tham số tác động lẫn nhau, và mỗi lần đánh giá cần chạy mô phỏng với nhiều seed. Chỉnh tay phụ thuộc vào kinh nghiệm người làm và khó tái lập. Tối ưu hộp đen tìm được giá trị số, nhưng không đề xuất được cấu trúc luật mới (cần luật nào, áp dụng cho nhóm tài xế nào) và không giải thích vì sao bộ tham số này tốt hơn bộ khác. Câu hỏi của luận văn là: làm thế nào xây dựng và tinh chỉnh bộ luật hành vi tài xế cho mô phỏng hỗn hợp một cách có hệ thống, có kiểm soát và tái lập được.

\section{Mục tiêu và phạm vi đề tài}
\label{section:1.2}

\subsection{Mục tiêu}

Mục tiêu của luận văn là mô phỏng hành vi tham gia giao thông tại nút Lê Văn Lương -- Hoàng Đạo Thúy giống thực tế nhất có thể. Để đạt mục tiêu này, một mô hình ngôn ngữ lớn (LLM) chạy cục bộ sinh và tinh chỉnh lặp tập luật hành vi tài xế cho SUMO. Luận văn có năm mục tiêu cụ thể:
\begin{enumerate}[label=(\roman*)]
    \item dựng mạng mô phỏng tái lập được cho nút giao, chỉ từ script và các tệp vá cấu hình;
    \item định nghĩa dạng biểu diễn luật gồm tham số kiểu phương tiện và luật có điều kiện thực thi qua TraCI, kèm bộ kiểm tra hợp lệ và bộ biên dịch;
    \item xây dựng bộ điều khiển thực thi luật trong khi mô phỏng, ghi lại mọi lần can thiệp;
    \item xây dựng hệ thống chỉ số để so sánh mô phỏng với dữ liệu quan sát;
    \item so sánh bộ luật do LLM sinh với SUMO mặc định, bộ luật thủ công và Optuna, trên cùng ngân sách số lần chạy; tách hai vai trò của LLM là chỉnh tham số trên cấu trúc cố định và sinh cấu trúc luật.
\end{enumerate}

\subsection{Phạm vi}

\textbf{Không gian và thời gian.} Luận văn chỉ xét một nút giao có đèn, bốn nhánh: Lê Văn Lương -- Hoàng Đạo Thúy (quận Thanh Xuân, Hà Nội), trong cao điểm sáng 07:30--09:00. Dòng xe gồm ô tô và xe máy. Người đi bộ, xe buýt và xe tải không được mô hình hóa.

\textbf{Kịch bản hành vi.} Kịch bản mô phỏng gồm 11 hành vi, chia thành ba nhóm (Bảng~\ref{table:1.1}). Rẽ trái trái phép chờ giữa nút là hành vi (10); nó là một trong 11 hành vi, cùng được mô hình hóa và đánh giá như các hành vi còn lại. Hành vi (8), đi ngược chiều, đã qua thí nghiệm khả thi nhưng được hoãn, nên chưa có trong mô hình.

\begin{table}[ht]
\centering
\caption{Mười một hành vi trong kịch bản mô phỏng.}
\label{table:1.1}
\begin{tabular}{|c|p{3.2cm}|p{8cm}|}
\hline
\textbf{STT} & \textbf{Nhóm} & \textbf{Hành vi} \\
\hline
1 & \multirow{3}{3.2cm}{Không tuân thủ tín hiệu} & Ngưỡng kiên nhẫn và vượt đèn đỏ \\
2 & & Áp lực bầy đàn: vượt đèn đỏ theo đoàn \\
3 & & Xuất phát sớm trước khi đèn xanh \\
\hline
4 & \multirow{5}{3.2cm}{Tìm khoảng trống khi ùn} & Khoảng cách ngang nhỏ, lấn chiếm và luồn lách \\
5 & & Đi sai làn để né đèn \\
6 & & Xe đi thẳng chiếm làn rẽ phải liên tục \\
7 & & Xe máy đi trên vỉa hè để vượt hàng chờ \\
8 & & Đi ngược chiều ở đoạn không có dải phân cách (hoãn) \\
\hline
9 & \multirow{3}{3.2cm}{Tương tác trong nút} & Chấp nhận khoảng trống nhỏ, tạt đầu \\
10 & & Rẽ trái trái phép: đi vào giữa nút, dừng chờ, rẽ khi pha đi thẳng của hướng đối diện kết thúc \\
11 & & Xe máy đi chậm giữa làn \\
\hline
\end{tabular}
\end{table}

\textbf{Điều LLM được và không được làm.} LLM sinh luật qua hai kênh: tham số kiểu phương tiện, và luật TraCI có điều kiện. LLM cũng quyết định tỷ trọng các kiểu tài xế. LLM không sinh hay sửa mạng lưới đường. Mạng được giữ cố định, ngoài các bản vá tĩnh chỉ dành cho lớp xe vi phạm: kết nối rẽ trái trái phép và làn vỉa hè, những hành vi cần không gian riêng.

\textbf{Ngoài phạm vi.} Ô tô đỗ dưới lòng đường và mô hình chọn lộ trình không được cài đặt.

\textbf{Dữ liệu.} Luận văn dự kiến dùng dữ liệu quan sát trích từ video quay tại nút trong nhiều ngày; xử lý video nằm ngoài phạm vi. Tại thời điểm viết, chưa có dữ liệu video. Các thông số hiện trường chưa đo được để trống có chủ đích. Quy trình được phát triển và kiểm thử trên dữ liệu đại diện (proxy), xây dựng từ tài liệu về các nút giao khác. Vì vậy kết quả trong luận văn không dùng để kết luận mô phỏng giống giao thông thực đến mức nào.

\section{Định hướng giải pháp}
\label{section:1.3}

Luận văn tách việc \emph{sinh luật} khỏi việc \emph{thực thi luật} (Hình~\ref{fig:1.1}). Ở mỗi vòng hiệu chỉnh, LLM nhận mô tả nút giao, từ điển các tham số được phép chỉnh, và phản hồi định lượng của vòng trước. LLM trả về một đặc tả luật dạng JSON. Đặc tả được kiểm tra hợp lệ rồi biên dịch thành hai thứ: tham số tĩnh của các kiểu phương tiện mà SUMO nạp khi khởi động, và các luật có điều kiện mà bộ điều khiển TraCI đánh giá ở mỗi bước mô phỏng. LLM không điều khiển xe theo thời gian thực. Kết quả mô phỏng được chuyển về cùng lược đồ với dữ liệu quan sát để tính chỉ số, và các chỉ số tạo thành phản hồi cho vòng sau. Phương pháp chi tiết nằm ở Chương~\ref{chapter:Methodology}.

\begin{figure}[ht]
\centering
\begin{tikzpicture}[
    node distance=8mm and 10mm,
    box/.style={draw, rounded corners=2pt, align=center, minimum height=11mm, text width=2.7cm, font=\small},
    llm/.style={box, fill=gray!20, very thick},
    arr/.style={-{Stealth[length=2.4mm]}, thick}
]
\node[llm] (llm) {LLM cục bộ:\\sinh đặc tả luật};
\node[box, right=of llm] (val) {1. Kiểm tra\\hợp lệ};
\node[box, right=of val] (cmp) {2. Biên dịch};
\node[box, below=of cmp] (sim) {3. Chạy SUMO\\và bộ điều khiển TraCI};
\node[box, left=of sim] (met) {4. Tính chỉ số\\M1--M10};
\node[box, below=of llm] (fb) {5. Phản hồi\\cho vòng sau};
\draw[arr] (llm) -- (val);
\draw[arr] (val) -- (cmp);
\draw[arr] (cmp) -- (sim);
\draw[arr] (sim) -- (met);
\draw[arr] (met) -- (fb);
\draw[arr] (fb) -- (llm);
\begin{pgfonlayer}{background}
\node[draw, dashed, gray, rounded corners=3pt, inner sep=4mm, fit=(val)(cmp)(sim)(met), label={[font=\scriptsize, gray]above:mã xác định, không gọi LLM}] {};
\end{pgfonlayer}
\end{tikzpicture}
\caption{Một vòng hiệu chỉnh. LLM chỉ chạy ở bước sinh đặc tả; các bước 1--4 là mã xác định. Bước 2 xuất hai tệp: \texttt{vtypes.add.xml} (tham số kiểu phương tiện) và \texttt{traci\_rules.json} (luật có điều kiện).}
\label{fig:1.1}
\end{figure}

Luận văn có bốn đóng góp:
\begin{itemize}
    \item mạng mô phỏng và quy trình chạy tái lập được cho nút giao, trong đó mỗi lần chạy ghi một tệp \texttt{manifest.json};
    \item dạng biểu diễn luật hành vi có miền tham số kiểm chứng theo tài liệu SUMO, cùng bộ kiểm tra hợp lệ, bộ biên dịch và bộ điều khiển TraCI;
    \item thí nghiệm khả thi cho các hành vi cần không gian riêng, trong đó có dừng chờ giữa nút;
    \item khung so sánh công bằng giữa luật do LLM sinh và các phương án đối sánh, trên cùng không gian tham số và cùng ngân sách mô phỏng.
\end{itemize}
Hệ thống chỉ số, vòng lặp LLM và phương án Optuna đã được cài đặt. Các chiến dịch hiệu chỉnh chưa xong, và đánh giá cuối vẫn là phần dự kiến.

\section{Bố cục luận văn}
\label{section:1.4}

% Lưu ý: tên chương lấy từ DoAn.tex; cập nhật khi đổi tên chương.
\begin{itemize}
    \item Chương~\ref{chapter:Related_works} khảo sát các nghiên cứu liên quan, nêu chức năng và yêu cầu của hệ thống.
    \item Chương~\ref{chapter:Methodology} trình bày phương pháp: kịch bản, mạng mô phỏng, biểu diễn luật, kiểm tra hợp lệ, bộ điều khiển, dữ liệu mục tiêu, chỉ số, quy trình hiệu chỉnh và các phương án đối sánh.
    \item Chương~\ref{chapter:Experiment} mô tả phần cài đặt và các kết quả thực nghiệm đã có, kèm điều kiện chạy của từng kết quả.
    \item Chương~\ref{chapter:SolutionAndContribution} tổng hợp các giải pháp kỹ thuật và đóng góp chính.
    \item Chương~\ref{chapter:conclusion} tóm tắt kết quả, hạn chế và hướng phát triển.
\end{itemize}

\end{document}
